% Generated draft: requires mathematical and source review.
\documentclass[11pt]{article}
\usepackage[margin=25mm]{geometry}
\usepackage{fontspec}
\setmainfont{Latin Modern Roman}
\usepackage{amsmath,amssymb,mathtools,unicode-math}
\setmathfont{Latin Modern Math}
\usepackage{enumitem,needspace,xurl,hyperref}
\setlength{\parindent}{0pt}
\setlength{\parskip}{5pt}
\setlength{\emergencystretch}{3em}
\newcommand{\N}{\mathbb{N}}
\newcommand{\Z}{\mathbb{Z}}
\newcommand{\Q}{\mathbb{Q}}
\newcommand{\R}{\mathbb{R}}
\newcommand{\C}{\mathbb{C}}
\newcommand{\F}{\mathbb{F}}
\newcommand{\A}{\mathbb{A}}
\newcommand{\PP}{\mathbb{P}}
\newcommand{\E}{\mathbb{E}}
\newcommand{\eps}{\varepsilon}
\newcommand{\dd}{\,\mathrm{d}}
\begin{document}
\section*{UnsolvedMath ID: 20001915}
\textbf{Source problem code:} catalogue entry
\subsection{Definitions and mathematical statement}
The source asks whether the union of closed geodesics is always dense in the manifold. The aid uses a connected closed smooth positive-dimensional Riemannian manifold \((M,g)\) and nonconstant closed geodesics. Allowing constant geodesics would make the question vacuous. This is base-space density, weaker than density of periodic vectors in the whole unit tangent bundle.

Let \(SM\) have its Sasaki metric, \(\pi:SM\to M\) be projection, and \(P_g\) the periodic vectors of the unit-speed geodesic flow. Write \(E_g=\overline{P_g}\) and
\[
 \Delta_g(x)=\operatorname{dist}_{SM}(S_xM,E_g).
\]
The metric choice matters to the numerical Lipschitz assertion; the source does not specify it, so Sasaki distance is explicit here. The aid reports
\[
 \overline{\bigcup_{c\text{ nonconstant closed geodesic}}c(\R)}
       =\pi(E_g)=\{x:\Delta_g(x)=0\},
 \qquad |\Delta_g(x)-\Delta_g(y)|\le d_g(x,y).
\]
Compactness/properness of the unit-fiber projection is used in identifying the projection of the closure; it is not an automatic projection rule on arbitrary noncompact bundles. Hence the original density question becomes whether \(\Delta_g\equiv0\), with a positive gap specifying an open base region missed by all closed geodesics. This reformulation does not prove gap vanishing.

The candidate compactness lemma uses smooth metrics \(g_j\to g\) in \(C^1\) on the fixed compact manifold. Suppose unit-speed \(g_j\)-closed geodesics \(c_j\) have periods \(0<T_j\le L\) and meet a fixed compact set \(K\subset U\), where \(U\) is a prescribed open base region. The aid reports that, after shifting the starting point to the hit and passing to a subsequence, they converge to a nonconstant closed \(g\)-geodesic meeting \(K\), with a positive limiting period. Compact convergence supplies the lower-period bound needed to rule out collapse to a constant loop. No nondegeneracy or continuation theorem for individual closed geodesics is assumed.

The compact-inside-open condition is material: if hits are only in \(U\), a subsequential limit can land on its boundary. For a density argument one can test a smaller region with closure inside \(U\). Accordingly density passes from a nearby-metric family to \(g\) if, in every nonempty target open set, such locally hitting geodesics exist with a uniform period bound. Nearby density without these bounds is not sufficient; periods may escape to infinity. The fiber-gap and period-escape descriptions identify limitations of this perturbative route, not a general proof or counterexample.

The aid reports the Anosov geodesic-flow case as a corollary of periodic-orbit density in that compact setting. It states the universal closed-Riemannian density problem remains open. All geometric and status assertions here are attributed unreviewed research-aid claims.
\subsection{Short English statement}
The aid rewrites base density as vanishing of a Sasaki fiber-gap function and gives a bounded-period compactness test for passing local closed-geodesic hits from nearby metrics to a limit. The Anosov case is reported covered, but the universal question and possible period escape remain unresolved.
\subsection{Sources}
\begin{itemize}
\item Local UnsolvedMath registry, permanent ID \(20001915\), source code \texttt{catalogue entry}; \texttt{aim-geometry-notes.json}, \texttt{notes[252]}, tag \texttt{question}.
\item Keith Burns and Vladimir S. Matveev's AIM ``Geodesics'' problem list, recorded item 8.5, Question 8.5.1: \url{https://aimath.org/pastworkshops/geodesicsproblems.pdf}. Source type PDF; statement status \texttt{exact}. Only the supplied record was used; no PDF was accessed or compiled.
\item No authoritative full bibliography, explicit metric convention, theorem identifier, or complete proof artifacts are supplied for the fiber-gap, compactness, or Anosov corollary. Original remarks and literature fields are empty.
\item Research attempt 1: \texttt{partial\_result}, status at run \texttt{open}, rigor label \texttt{proved}, novelty confidence low; registry status \texttt{partially\_solved}, classification \texttt{PARTIAL-PROGRESS}. Created and updated 14 August 2026. Machine-generated aid, not peer reviewed; independently verify before citation.
\item Ulam AI / UnsolvedMath Contributors; CC BY 4.0.
\end{itemize}
\end{document}
